\subsection{Finite-Dimensional Splitting and Isomorphism Tests}
\begin{definition} \label{def:codimension-subspace} \hlt{Codimension}\\
If $\mathcal{U}$ is a subspace of $\mathcal{V}$ and $\dim(\mathcal{V}/\mathcal{U})<\infty$, define $\operatorname{codim}_{\mathcal{V}}(\mathcal{U})\defeq\dim(\mathcal{V}/\mathcal{U})$.\\
Thus, if $\mathcal{V}$ is finite-dimensional, Theorem \ref{thm:dimension-of-quotient} gives $\operatorname{codim}_{\mathcal{V}}(\mathcal{U})=\dim(\mathcal{V})-\dim(\mathcal{U})$.\\
If $\mathcal{V}=\mathcal{U}\oplus\mathcal{W}$, Proposition \ref{prop:quotient-by-complement-is-complement} gives $\operatorname{codim}_{\mathcal{V}}(\mathcal{U})=\dim(\mathcal{W})$.
\end{definition}

\begin{perspective} \label{per:quotients-split-over-fields} \hlt{Quotients Split Over Fields}\\
In finite-dimensional vector spaces, choosing a complement $\mathcal{V}=\mathcal{U}\oplus\mathcal{W}$ identifies $\mathcal{V}/\mathcal{U}$ with $\mathcal{W}$. The quotient map collapses $\mathcal{U}$ and keeps the complementary directions.
\end{perspective}

\begin{proposition} \label{prop:finite-dimensional-one-sided-inverses-split} \hlt{One-Sided Inverses Split in Finite Dimensions}\\
Let $\mathcal{V}$ and $\mathcal{W}$ be finite-dimensional vector spaces and let $T\in\LL(\mathcal{V},\mathcal{W})$.
\begin{enumerate}[label=(\roman*)]
\setlength{\itemsep}{0pt}
\item If $T$ is injective, then there exists $L\in\LL(\mathcal{W},\mathcal{V})$ such that $LT=I_{\mathcal{V}}$.
\item If $T$ is surjective, then there exists $R\in\LL(\mathcal{W},\mathcal{V})$ such that $TR=I_{\mathcal{W}}$.
\end{enumerate}
\begin{proof}
\begin{enumerate}[label=(\roman*)]
\setlength{\itemsep}{0pt}
\item Choose $\mathcal{H}\subseteq\mathcal{W}$ with $\mathcal{W}=\operatorname{range}(T)\oplus\mathcal{H}$. For $\mathbf{w}=T\mathbf{v}+\mathbf{h}$, define $L\mathbf{w}=\mathbf{v}$. This is well-defined since the direct-sum decomposition is unique and $T$ is injective. It is linear, and $LT\mathbf{v}=\mathbf{v}$ for every $\mathbf{v}\in\mathcal{V}$.
\item Choose $\mathcal{H}\subseteq\mathcal{V}$ with $\mathcal{V}=\ker(T)\oplus\mathcal{H}$. Then $T|_{\mathcal{H}}:\mathcal{H}\to\mathcal{W}$ is injective since $\mathcal{H}\cap\ker(T)=\{\mathbf{0}\}$, and surjective since each $\mathbf{v}=\mathbf{k}+\mathbf{h}$ satisfies $T\mathbf{v}=T\mathbf{h}$. Let $R=(T|_{\mathcal{H}})^{-1}$, viewed as a map $\mathcal{W}\to\mathcal{V}$.\\
Then $R$ is linear and $TR=I_{\mathcal{W}}$.
\end{enumerate}
\hlt{Structural Commentary}: Complements supply the missing directions on which a one-sided inverse is chosen.\\
\hlt{Key Insight}: Finite-dimensional vector spaces split kernels and ranges.
\end{proof}
\end{proposition}

\begin{proposition} \label{prop:isomorphisms-preserve-bases} \hlt{Isomorphisms Preserve Bases}\\
Let $T:\mathcal{V}\to\mathcal{W}$ be an isomorphism and let $\beta=(\mathbf{v}_1,\ldots,\mathbf{v}_n)$ be a finite list in $\mathcal{V}$. Set $T\beta=(T\mathbf{v}_1,\ldots,T\mathbf{v}_n)$.\\
The following are equivalent:
\begin{enumerate}[label=(\roman*)]
\setlength{\itemsep}{0pt}
\item $\beta$ is a basis of $\mathcal{V}$.
\item $T\beta$ is a basis of $\mathcal{W}$.
\end{enumerate}
\begin{proof}
\begin{enumerate}[label=(\roman*).]
\setlength{\itemsep}{0pt}
\item[$(i)$]$\Rightarrow (ii)$ If $\mathbf{w}\in\mathcal{W}$, choose $\mathbf{v}\in\mathcal{V}$ with $\mathbf{w}=T\mathbf{v}$ and write $\mathbf{v}=\sum_j a_j\mathbf{v}_j$. Then $\mathbf{w}=\sum_j a_jT\mathbf{v}_j$, so $T\beta$ spans $\mathcal{W}$. If $\sum_j a_jT\mathbf{v}_j=\mathbf{0}$, applying $T^{-1}$ gives $\sum_j a_j\mathbf{v}_j=\mathbf{0}$; hence all $a_j=0$. Thus $T\beta$ is a basis.
\item[$(ii)$]$\Rightarrow (i)$
Assume $T\beta$ is a basis of $\mathcal{W}$. If $\mathbf{v}\in\mathcal{V}$, write $T\mathbf{v}=\sum_j a_jT\mathbf{v}_j$; applying $T^{-1}$ gives $\mathbf{v}=\sum_j a_j\mathbf{v}_j$, so $\beta$ spans $\mathcal{V}$. If $\sum_j a_j\mathbf{v}_j=\mathbf{0}$, applying $T$ gives $\sum_j a_jT\mathbf{v}_j=\mathbf{0}$; hence all $a_j=0$. Thus $\beta$ is linearly independent, and therefore a basis of $\mathcal{V}$.
\end{enumerate}
\hlt{Structural Commentary}: Isomorphisms preserve exactly the spanning and independence data that define a basis.\\
\hlt{Key Insight}: An isomorphism changes the ambient space, not the linear relations among vectors.
\end{proof}
\end{proposition}

\begin{theorem} \label{thm:finite-dimensional-operator-invertible-equivalences} \hlt{Finite-Dimensional Invertibility Criteria}\\
If $\mathcal{V}$ is finite-dimensional and $T \in \LL(\mathcal{V})$, then the following are equivalent:
\begin{enumerate}[label=\roman*.]
\setlength{\itemsep}{0pt}
\item $T$ is invertible;
\item $T$ is injective;
\item $T$ is surjective.
\end{enumerate}
\begin{proof}
\begin{enumerate}[label=(\roman)*.]
\setlength{\itemsep}{0pt}
\item[$(i)$] $\Rightarrow (ii)$ follows from Proposition \ref{prop:linear-map-invertible-iff-bijective}.
\item[$(ii)$] $\Rightarrow (iii)$: If $T$ is injective, then $\ker(T)=\{\mathbf{0}\}$.\\
Rank-nullity gives $\dim(\operatorname{range}(T))=\dim(\mathcal{V})$, hence $\operatorname{range}(T)=\mathcal{V}$, so $T$ is surjective.
\item[$(iii)$] $\Rightarrow (i)$: If $T$ is surjective, rank-nullity gives $\dim(\ker(T))=0$, so $T$ is injective.\\
Thus $T$ is bijective, hence invertible by Proposition \ref{prop:linear-map-invertible-iff-bijective}.
\end{enumerate}
\hlt{Structural Commentary}: For operators on finite-dimensional spaces, kernel and range dimensions trade off exactly. Injectivity and surjectivity therefore force each other.\\
\hlt{Key Insight}: In equal finite dimensions, losing no information is the same as reaching the whole space.
\end{proof}
\end{theorem}

\begin{perspective} \label{per:kernel-obstruction-to-invertibility} \hlt{Kernel Obstruction to Invertibility}\\
By Theorem \ref{thm:finite-dimensional-operator-invertible-equivalences} and Proposition \ref{prop:injective-iff-kernel-zero}, a finite-dimensional operator is invertible exactly when its kernel is zero. Any nonzero $\mathbf{v}\in\ker T$ therefore blocks invertibility.
\end{perspective}

\begin{theorem} \label{thm:finite-dimensional-isomorphic-iff-same-dimension} \hlt{Isomorphism and Dimension}\\
Two finite-dimensional vector spaces are isomorphic if and only if they have the same dimension.
\begin{proof}
\begin{enumerate}[label=\roman*.]
\setlength{\itemsep}{0pt}
\item[$(\Rightarrow)$] If $T:\mathcal{V}\to\mathcal{W}$ is an isomorphism and $(\mathbf{v}_1,\ldots,\mathbf{v}_n)$ is a basis of $\mathcal{V}$, then Proposition \ref{prop:isomorphisms-preserve-bases} shows that $(T\mathbf{v}_1,\ldots,T\mathbf{v}_n)$ is a basis of $\mathcal{W}$.
\item[$(\Leftarrow)$] If both spaces have dimension $n$, choose bases $(\mathbf{v}_1,\ldots,\mathbf{v}_n)$ and $(\mathbf{w}_1,\ldots,\mathbf{w}_n)$.\\
By Proposition \ref{prop:linear-map-determined-by-basis}, there is a linear map $T$ with $T\mathbf{v}_j=\mathbf{w}_j$ for every $j$. Then $T(\sum a_j\mathbf{v}_j)=\sum a_j\mathbf{w}_j$, so $\ker(T)=\{\mathbf{0}\}$ and $\operatorname{range}(T)=\mathcal{W}$. Hence $T$ is an isomorphism.
\end{enumerate}
\hlt{Structural Commentary}: Dimension classifies finite-dimensional spaces up to isomorphism.\\
\hlt{Key Insight}: Matching bases give an isomorphism; any isomorphism carries bases to bases.
\end{proof}
\end{theorem}
